\begin{preface}
\addcontentsline{toc}{chapter}{Preface} %% adds this to the table of content
\markboth{}{} %% removes header from pages

This thesis has in a way been found through ``experimental'' mathematics, that is that especially in the final 6 months, we have used Macaulay2 to find.
\end{preface}

% \begin{preface}
% \addcontentsline{toc}{chapter}{Preface} %% adds this to the table of content
% \markboth{}{} %% removes header from pages
% The preface is a short and perhaps more personal description of the project, the purpose and your motivation for pursuing it. It should be short and sweet; and it is different from the acknowledgements. 

% This template has grown out of the work with seminars about scientificness and writing in mathematical subjects that was organized by the Department of Mathematics and the Science Library at the University of Oslo (UiO), and hosted by Karoline Moe, Ph.D., in the period from 2014 to 2024. 

% There has always been a tradition at the department that former students pass their \LaTeX-template on to the next generation of students. However, this practice did not ensure compliance neither with standard traditions in scientific writing and reference management, nor with subsequent design manuals at UiO and typography. Moreover, some of the commands and packages in \LaTeX\ used by students were outdated or in conflict, and there was clearly a need for more instruction. 

% This template tries to answer these issues. The main purpose of the template is to serve as a starting point for all students using \LaTeX\ to write their Master's thesis. First of all it is a \enquote{minimal} working example with the structure of a thesis. Secondly, We have included examples that most students frequently run into and that can take forever to figure out, even with all available tools, search engines and AI. Finally, we have tried to include a description or short text of the most essential genre-features or pit-falls for each chapter. 

% We sincerely hope that this template can be of some help in your work, and feel free to send the the \LaTeX-group at the University of Oslo Library an e-mail on \hyperlink{latexguru@ub.uio.no}{latexguru@ub.uio.no} if you are stuck!
% \end{preface}